\chapter{系统实现与使用}

\section{公共协议与原语分别实现}
系统分为应用入口、组合协议和原语三层。应用入口提供命令行与浏览器操作；组合协议负责身份编码、混合封装、完整密文认证和状态管理；HIBE与IBS分别完成各自密码运算。图~\ref{fig:system-structure}给出模块关系。

\begin{figure}[H]
\centering
\begin{tikzpicture}[node distance=8mm and 10mm]
 \node[flow,text width=43mm] (web) {浏览器交互\\配置、收发、结果展示};
 \node[flow,right=of web,text width=43mm] (cli) {命令行入口\\批处理、保存与恢复};
 \node[check,below=of web,text width=43mm] (protocol) {统一IB-ME协议\\身份、封装与认证};
 \node[check,right=of protocol,text width=43mm] (state) {规范对象与状态管理\\参数检查、原子保存};
 \node[flow,below=of protocol,text width=43mm] (hibe) {HIBE公开接口\\LATTE、Lai};
 \node[flow,right=of hibe,text width=43mm] (ibs) {IBS公开接口\\ML-DSA转换、SRGG、Sageloli};
 \draw[flowarrow](web)--(protocol);\draw[flowarrow](cli)--(state);
 \draw[flowarrow](protocol)--(state);
 \draw[flowarrow](protocol)--(hibe);\draw[flowarrow](state)--(ibs);
 \draw[flowarrow](protocol.south east)--(ibs.north west);
\end{tikzpicture}
\caption{应用通过统一协议调用独立HIBE与IBS组件}
\label{fig:system-structure}
\end{figure}

HIBE接口包括建立、一级提取、二级委派、加密和解密；IBS接口包括建立、身份发钥、签名和验证。公共层通过规范对象交互，内部多项式、矩阵和陷门保留在原语实现中。更换组件时，先检查两层身份能力、内部身份宽度、签名输入上限、签名尺寸与HIBE载荷容量，再执行密码运算。

除三组代表外，系统还支持BNN/Falcon-512和Luo模块格IBS~\cite{falcon2020,luo2024ibs}。两种HIBE与五类IBS共十种配对的功能测试，验证了同一协议在不同组件下的组合能力。报告的主要性能分析集中于A、B、C。

\section{固定发钥、保存与重启恢复}
发钥机构为同一身份保存已经签发的结果，避免同一用户在后续调用中无意获得另一份密钥。一次发钥操作包含三个阶段：记录开始状态、执行密码运算、提交完成状态。应用在本地锁保护下先保存主状态，再原子发布用户密钥。

若进程在发钥中途终止，状态保留未完成标记，由恢复过程处理。加载公钥、主状态和用户钥时，系统检查类型、版本、参数、公钥上下文与身份绑定，并执行相应原语方程检查。命令行能够从持久对象重新启动，继续解密已有密文。本地状态文件作为秘密数据保存，部署环境负责访问控制与备份。

\section{一次完整的交互操作}
浏览器界面展示接收者、发送者、期望发送者、消息和运行结果。使用者首先选择配置并建立系统，随后按以下顺序完成一次收发：
\begin{enumerate}
 \item 输入双方身份，为发送者与接收者生成密钥。
 \item 输入消息并指定接收身份，生成完整密文。
 \item 使用接收密钥，输入期望发送者，执行解密。
 \item 查看恢复消息、密文长度及本次耗时。
 \item 在同一密文上切换身份错配或篡改模式，查看拒绝结果。
\end{enumerate}
图~\ref{fig:web-demo}来自A组实际运行：发送者Alice面向Bob加密36字节UTF-8消息，正确匹配时恢复原文。同一密文还用于三种身份错配、字节修改和截断测试。

\begin{figure}[H]
\centering
\includegraphics[width=.98\textwidth]{../../deliveries/2026-10-20-web-redesign/desktop.png}
\caption{A组实际收发界面：指定双方身份、加密消息与匹配解密}
\label{fig:web-demo}
\end{figure}

A、B两组均完成实际浏览器收发和拒绝测试。B的发钥和签名时间较长，界面持续显示任务阶段，页面刷新后接回同一个任务，避免重复启动密码操作。公开参数以精确整数传输与导出，支持超过浏览器普通整数范围的参数。

浏览器测试中，两组分别在同一条密文上执行正确解密、三种身份错配、翻转1 bit和截短1字节，结果均符合预期。A、B该次密文长度为142,095字节和149,255字节。

网页刷新保留的是仍在运行的服务会话；服务重启后的恢复由命令行加载持久对象完成。C的完整长流程、恢复和接收会话已有命令行与公开接口实验，浏览器入口用于配置和环境检查。

\section{减少运行中的重复工作}
\subsection{证书验证与采样准备缓存}
BNN签名包含机构证书和用户消息签名。对同一身份的连续消息，可以缓存“该机构证书验证成功”的公开结果，逐条消息仍验证用户签名。缓存绑定完整证书、公钥与身份，采用有界LRU管理，默认保留64项。

SRGG签名中的部分数值准备取决于固定签名钥与数值选项。实现可保留这些计算计划，默认最多2项；后续签名重新采样随机量。缓存节省的是重复准备的工作，其完整加密收益由第5章的配对实验说明。

\subsection{部署文件的位置与读取开销}
原生组件在建立适配器时校验源码、接口和二进制。完整流程剖析发现，跨Windows挂载目录反复读取文件会增加时间。系统保留这些检查，将固定安装包与运行文件放在WSL的Linux文件系统中，减少挂载路径的访问成本。

第5章用同一WSL运行时、相同密码参数及同一安装包比较文件放置位置。该结果帮助选择本机运行方式，与前述采样、编码和会话方法共同构成完整系统的实现优化。

\section{独立安装与运行检查}
另在Xeon Silver 4216主机的NixOS、Sage 10.9／Python 3.14.7中，以独立安装包复现三组功能，OMP与BLAS均为1。A、B各完成25条操作，C完成匹配解密、内部签名拒绝与恢复。C最初900秒阶段上限导致超时，调整为3600秒后完成，三个阶段耗时约21至27分钟。完整运行记录保留在支撑材料。
